% Propietario conservado: /Users/ruben/Documents/New project/output/EXTREMA_MEDIA_RAZON_RECIPROCIDAD_ES_EN_20260918/ARTICULO/ES/source/sections/03_pantallas_coordinadas.tex
% Procedencia: integral 2249, cap. 25, pp. 619--640, especialmente §25.13.
% Requiere la incorporación material de exc:k-direccion, exc:leech y el núcleo común.
\section{Pantallas dimensionales y reducción coordinada del estado}
\label{x_reciprocidad:iv:pantallas}

Los rangos que aparecen en la prolongación excepcional y en la
compactificación pertenecen a objetos diferentes. Para relacionarlos
se conserva su presentación: generadores, relaciones, producto
interno y dirección seleccionada. La dimensión es entonces una
consecuencia del módulo de canales y de sus relaciones. Su
transporte se formula mediante mapas que actúan sobre ese módulo,
además de actuar sobre el estado que lo determina.

\subsection{Canales, relaciones y dimensión efectiva}

Sea \(\mathcal S_n\) el conjunto de estados supervivientes de
profundidad \(n\), con restricciones compatibles
\(r_{m,n}:\mathcal S_m\to\mathcal S_n\), \(m\geq n\).
El dato de frontera \(\beta_n(x)\) conserva la información
necesaria para decidir la compatibilidad de una prolongación:
hoja, orientación, cociente, memoria y clase de incidencia según
el lector considerado. Las restricciones satisfacen
\[
 r_{\ell,n}=r_{m,n}r_{\ell,m},\qquad
 \beta_n r_{m,n}=b_{m,n}\beta_m.
\]
Estas igualdades son las condiciones de naturalidad de la
presentación y mantienen los mismos datos al restringir por
etapas o directamente.

\begin{definition}[Pantalla de canales]
\label{x_reciprocidad:iv:def:pantalla}
Fijado un cuerpo \(\mathbb K_n\), sea \(M_n(x)\) el espacio
generado por los canales de prolongación y \(\mathcal R_n(x)\)
el subespacio de relaciones. La pantalla es
\[
 \Sigma_n(x)=\bigl(\mathbb K_n,M_n(x),\beta_n(x),
                    \mathcal R_n(x)\bigr),
\]
y su dimensión efectiva es
\begin{equation}
 d_{\Sigma_n}(x)=\dim_{\mathbb K_n}
                    \bigl(M_n(x)/\mathcal R_n(x)\bigr).
 \label{x_reciprocidad:iv:eq:dimension-pantalla}
\end{equation}
Si \(M_n(x)\cong\mathbb K_n^{g_n}\) y las columnas de
\(R_n(x)\) generan las relaciones, entonces
\(d_{\Sigma_n}(x)=g_n-\operatorname{rank}R_n(x)\).
\end{definition}

\begin{proposition}[Invariancia de presentación y transporte]
\label{x_reciprocidad:iv:prop:pantalla-transporte}
La dimensión \eqref{x_reciprocidad:iv:eq:dimension-pantalla} se conserva
bajo cambios de bases invertibles. Para dos pantallas sobre un mismo cuerpo
\(\mathbb K_n=\mathbb K_m=\mathbb K\), si la aplicación lineal
\(F_\gamma:M_n(x)\to M_m(T_\gamma x)\) envía
\(\mathcal R_n(x)\) en \(\mathcal R_m(T_\gamma x)\), induce
de manera única una aplicación lineal entre los cocientes de pantalla.
\end{proposition}
\begin{proof}
Un cambio de generadores y de relaciones transforma \(R\) en
\(URV\), con \(U,V\) invertibles, y conserva su rango. Para
el descenso, definimos
\(\overline F_\gamma([v])=[F_\gamma(v)]\). Si dos
representantes difieren en una relación, la linealidad implica que sus imágenes difieren
en una relación de la pantalla terminal; la aplicación está
bien definida. La proyección al cociente es sobreyectiva y
garantiza la unicidad.
\end{proof}

El cambio dimensional durante el transporte queda determinado
por
\[
 \Delta d(\gamma;x)=d_{\Sigma_m}(T_\gamma x)
                         -d_{\Sigma_n}(x).
\]
Esta expresión conserva el estado y su presentación de canales.
Un rango aislado no permite reconstruir ni las relaciones ni
el mapa que lo produjo.

\subsection{Perforación ternaria y pantalla de síndromes}

El código \(C_W=[12,6,6]_3\), construido a partir de la
incidencia excepcional, produce otra pantalla mediante
perforación de una coordenada. Denotemos por
\(p:C_W\to\FF_3^{11}\) ese mapa y por
\(G_{11}=p(C_W)\) su imagen.

\begin{proposition}[Código perforado y partición perfecta]
\label{x_reciprocidad:iv:prop:codigo-once}
El mapa \(p\) es inyectivo. El código \(G_{11}\) tiene
parámetros \([11,6,5]_3\), y sus bolas de Hamming de radio
\(2\) particionan \(\FF_3^{11}\).
\end{proposition}
\begin{proof}
Un vector del núcleo de \(p\) tendría soporte contenido en
la única coordenada retirada. La distancia mínima \(6\)
de \(C_W\) excluye todo vector no nulo de esa forma.
Así, la dimensión sigue siendo \(6\). Perforar disminuye
el peso a lo sumo en una unidad, luego la distancia mínima
de la imagen es al menos \(5\). Las bolas de radio \(2\)
son por ello disjuntas. Cada una contiene
\[
 1+2\binom{11}{1}+2^2\binom{11}{2}
 =1+22+220=243
\]
vectores. Como \(|G_{11}|=3^6\) y
\(3^6\cdot243=3^{11}\), las bolas cubren el ambiente.
Para comprobar que la distancia es exactamente \(5\),
tómese un vector de peso \(3\). Su bola decodificadora
no puede tener centro cero y posee un centro de peso a lo
sumo \(3+2=5\); dicho centro debe tener peso al menos
\(5\). Se obtiene así una palabra de peso \(5\).
\end{proof}

El entero \(243=3^{11-6}\) es aquí el número de clases
del cociente de síndromes \(\FF_3^{11}/G_{11}\). Su
función se fija por la perforación y la distancia del
código; la igualdad de cardinal con otra fibra no sustituye
el mapa entre ambas. Esta pantalla combinatoria y la
siguiente pantalla real conservan sus operaciones propias.

\subsection{La dirección dodecafásica y la reducción de rango}

La construcción de \ref{x_reciprocidad:exc:k-direccion} parte del registro
\(K\) ya obtenido por APP--TRIT--TPK y de la carta de
representación declarada sobre \(\RR^{12}\). En esa carta,
\begin{equation}
 V_{11}=\mathbf1^\perp,\qquad
 P_{11}=I_{12}-\frac1{12}J_{12},\qquad
 u_K=P_3P_{11}K.
 \label{x_reciprocidad:iv:eq:p11-uk}
\end{equation}
Aquí \(J_{12}\) es la matriz de unos y \(P_3\) es el
proyector isótipo especificado por la acción de \(A_5\).
La suma finita que define \(P_3\), sus representantes y la
generación de \(K\) se conservan en la sección citada de
este mismo artículo. En particular,
\[
 \|u_K\|^2=\frac{6638585+2275584\sqrt5}{20}>0.
\]
La norma positiva define \(\ell_K=\RR u_K\), y el
proyector
\begin{equation}
 P_{10}=P_{11}-\frac{u_Ku_K^{\mathsf T}}{\|u_K\|^2}
 \label{x_reciprocidad:iv:eq:p10}
\end{equation}
tiene imagen \(V_{10}=V_{11}\cap u_K^\perp\).

\begin{proposition}[Dos reducciones ortogonales]
\label{x_reciprocidad:iv:prop:12-11-10}
Se cumplen
\[
 \RR^{12}=\RR\mathbf1\mathbin{\oplus^\perp}V_{11},\qquad
 V_{11}=\ell_K\mathbin{\oplus^\perp}V_{10},\qquad
 \dim V_{10}=10.
\]
Además, \(P_{10}P_{11}=P_{10}\), \(P_{10}u_K=0\), y
\(P_{10}\) es un proyector ortogonal.
\end{proposition}
\begin{proof}
La proposición \ref{x_reciprocidad:prop:k-reduccion-diez}, demostrada antes de
construir las pantallas, establece la idempotencia, simetría, rango y
las dos identidades operatorias para este mismo \(u_K\) y este mismo
\(P_{10}\). La primera descomposición se obtiene de
\(P_{11}=I-J_{12}/12\); la segunda es la descomposición ortogonal
de aquella proposición. Así se conserva la demostración única de la
reducción y se explicitan aquí sus dos etapas.
\end{proof}

La primera reducción elimina exclusivamente el modo uniforme.
La segunda elimina la dirección determinada por \(K\) en
la representación fijada. Si se cambia simultáneamente la carta
por una matriz ortogonal \(S\), el vector se transforma por
\(u_K\mapsto Su_K\) y el proyector por
\(P_{10}\mapsto SP_{10}S^{-1}\). Esta covariancia conserva
la operación geométrica, aunque cambien sus coordenadas.

\subsection{Pantalla reticular lorentziana}

Sea \(\Lambda=N_\alpha\) el retículo de Leech construido
en \ref{x_reciprocidad:exc:leech}. Se considera el plano hiperbólico integral
\(U=\ZZ e_+\oplus\ZZ e_-\), con
\[
 \langle e_+,e_+\rangle=\langle e_-,e_-\rangle=0,
 \qquad\langle e_+,e_-\rangle=1.
\]
El retículo \(L_{26}=\Lambda\oplus U\) es par y
unimodular, de firma \((25,1)\). Cada vector se escribe
\(y=\lambda+ae_++be_-\). Como
\(\langle y,e_+\rangle=b\),
\begin{equation}
 e_+^\perp=\Lambda\oplus\ZZ e_+,
 \qquad e_+^\perp/\ZZ e_+\cong\Lambda.
 \label{x_reciprocidad:iv:eq:26-24}
\end{equation}
El mapa del cociente es
\(q_{24}(\lambda+ae_+)=\lambda\); es sobreyectivo y tiene
núcleo \(\ZZ e_+\). La disminución de rango en dos procede
de una restricción de ortogonalidad seguida por un cociente
isotrópico. Es, por tanto, una operación de tipo diferente a
\eqref{x_reciprocidad:iv:eq:p10}.

\subsection{Morfismo conjunto de pantallas}

Con el dominio \(\mathscr D_0\) de \ref{x_reciprocidad:iv:def:dualidad},
sean
\[
 \mathscr G_K=\{(v_{11},P_{10}v_{11}):v_{11}\in V_{11}\},
\]
\[
 \mathscr S_{\mathrm{HMT}}=\RR^{12}\times\mathscr D_0\times e_+^\perp,
 \qquad
 \mathscr S_M^{\mathrm{alg}}=\mathscr G_K\times\mathscr D_0\times\Lambda.
\]
\begin{definition}[Mapa algebraico de pantallas coordinadas]
\label{x_reciprocidad:iv:def:rm-alg}
Se define
\begin{equation}
 \mathcal R_M^{\mathrm{alg}}(v,d,y)
 =\bigl((P_{11}v,P_{10}v),d,q_{24}(y)\bigr).
 \label{x_reciprocidad:iv:eq:rm-alg}
\end{equation}
\end{definition}

\begin{theorem}[Isomorfismo cocientado y entrelazamiento de dualidad]
\label{x_reciprocidad:iv:thm:pantallas-isomorfas}
\label{x_reciprocidad:thm:k-pantallas-coordinadas}
El mapa \eqref{x_reciprocidad:iv:eq:rm-alg} induce un isomorfismo
\[
 \overline{\mathcal R}_M^{\mathrm{alg}}:
 (\RR^{12}/\RR\mathbf1)\times\mathscr D_0
 \times(e_+^\perp/\ZZ e_+)
 \longrightarrow\mathscr S_M^{\mathrm{alg}}.
\]
Entrelaza las transformaciones que actúan por \(\mathsf T_0\)
en el factor \(\mathscr D_0\) y por la identidad en los
otros dos factores. La forma \(E_{R_0}(m,w)\) se conserva
bajo esa dualidad.
\end{theorem}
\begin{proof}
La aplicación \([v]\mapsto P_{11}v\) es un isomorfismo
de \(\RR^{12}/\RR\mathbf1\) sobre \(V_{11}\).
Como \(P_{10}P_{11}=P_{10}\), el mapa
\([v]\mapsto(P_{11}v,P_{10}v)\) identifica ese cociente
con \(\mathscr G_K\); su inversa toma la primera componente
y su clase. Por \eqref{x_reciprocidad:iv:eq:26-24}, \(q_{24}\) induce el
isomorfismo reticular del tercer factor. El segundo se
transporta mediante la identidad. El producto de los tres
mapas prueba el isomorfismo.

Los proyectores y \(q_{24}\) no modifican \(d\), mientras
la dualidad actúa únicamente en \(d\). Por ello las dos
composiciones son idénticas. La conservación de la forma
es el teorema \ref{x_reciprocidad:iv:thm:dualidad-escalar}.
\end{proof}

El grafo \(\mathscr G_K\) conserva la coordenada anterior
\(v_{11}\) junto con su proyección. El teorema demuestra
la equivalencia de las presentaciones cocientadas; la reducción
aislada \(V_{11}\to V_{10}\) conserva su núcleo
\(\ell_K\). Esta precisión establece qué información
retiene el isomorfismo y qué información elimina cada
proyección que lo compone.

\begin{corollary}[Pantallas sobre la familia elíptica]
El mismo morfismo se restringe a un isomorfismo
\[
 (\RR^{12}/\RR\mathbf1)\times\mathscr D_{\mathrm{el}}
 \times(e_+^\perp/\ZZ e_+)
 \ \cong\ \mathscr G_K\times\mathscr D_{\mathrm{el}}\times\Lambda,
\]
y conserva el entrelazamiento de dualidad.
\end{corollary}
\begin{proof}
El morfismo y su inversa conservan exactamente el segundo factor.
El corolario \ref{x_reciprocidad:iv:cor:radio-especializado} prueba que la dualidad
preserva \(\mathscr D_{\mathrm{el}}\). Por tanto se restringen
tanto la biyección como el cuadrado de entrelazamiento ya demostrado.
\end{proof}

\subsection{Coordinación con la realización excepcional}

Sobre el dominio \(\mathcal X_{\mathrm{HMT}}\), con fase, orientación,
cociente, frontera, incidencia, gradación y memoria, las lecturas
\(\varepsilon_{\mathrm{exc}}\) y \(\varepsilon_{\mathrm{scr}}\) extraen
los datos excepcionales y de pantalla. El corredor construye el álgebra
graduada \(\mathfrak V(x)\cong V^\natural\), con su producto de vértices
y sus automorfismos; \(\mathfrak S_M(x)\) reúne las pantallas, cocientes
y operadores anteriores. La asignación conjunta
\(x\mapsto(\mathfrak V(x),\mathfrak S_M(x))\) tiene, por tanto, un
álgebra como primer componente y una familia de espacios y mapas como segundo.
La acción del Monstruo pertenece al primero; la reducción dirigida por \(K\)
y la dualidad compacta, al segundo. Las dos lecturas fijan la procedencia
común y cada construcción conserva su tipo.
El isomorfismo de esta sección establece la correspondencia de
pantallas con sus factores y su prueba
explícitos.
